% Ad Quintum Fratrem 1.1 --- headnote
% ad-quintum-fratrem-01-01, late 60 BC or early 59 BC, written at Rome

\textit{Cicero to his brother Quintus, governor of Asia,
written at Rome at the end of 60 BC or the very beginning
of 59 BC. The longest of all the surviving letters and the
nearest thing in the corpus to a treatise: a forty-six-
section essay on provincial administration, addressed to
Quintus on the news that his year-long Asian command has
been extended for a third year (the announcement of the
extension is the letter's prompt, and Cicero takes blame
for it in \textsection 2: he had himself, while consul,
prevented the appointment of a successor in the previous
year). The letter is the first piece of the political-
philosophical body of Cicero's mature prose, written before
the dialogues of the late 50s and the philosophical works of
the 40s; it stands as their template.}

\textit{The argument moves through three large fields. The
first (\textsection 4--16) is the moral foundation: the
governor's own integrity, then the discipline of his
household, then the careful selection of those familiars ---
provincial Romans, Greeks, freedmen, slaves --- who could
trade on his name. ``Let your seal-ring be not as some
implement, but as you yourself are; not the agent of
another's will, but the witness of your own.'' Cicero
warns against trusting any provincial whose love for the
governor was not visible before the magistracy; against
the deceitful lightness of most contemporary Greeks
(``trained by long servitude to excessive flattery''),
without however excluding the few worthy of old Greece;
and against the use of slaves in any part of public
business.}

\textit{The second field (\textsection 17--26) is the
practice of justice. Severity is essential and ``must not
be varied by favour, but kept consistent''; but it must
be tempered by ``an easy bearing in hearing, mildness in
deciding, diligence in giving satisfaction and in
argument'' --- the model is the late Gaius Octavius (the
father of the future Augustus), in whose court ``the
first lictor was quiet, the herald silent, every man
spoke as long as he wished.'' The famous citation of
Xenophon's \textit{Cyropaedia} as the manual of just rule
that Africanus had at his elbow falls in \textsection 23.
The ruler's deepest principle (\textsection 24) is that
``everything must be referred to this end: that those who
shall be in their command be as happy as possible.'' The
list of Quintus's actual achievements --- the Mysian
brigandage put down, Samos and Halicarnassus rebuilt,
\textit{calumnia} (the false accusation that was the
``bitterest handmaid of praetorial greed'') driven out ---
fills \textsection 25.}

\textit{The third field (\textsection 32--36) is the
politically delicate one: the publicans. The Roman
equestrians who farmed the Asian taxes were Cicero's own
political base; they were also the chief instrument of
provincial misery. Cicero's prescription is the
characteristic both/and: protect both the publicans and
the Greeks, ``which seems the work of a divine kind of
virtue --- which is yours.'' Make the Greeks see that
without Roman command they would be tributary anyway,
and worse off; let the publicans see that without
Quintus's authority they cannot collect; bring them
together by the governor's own gravity and favour.}

\textit{The fourth and last field is the personal one. The
exception in the universal praise of Quintus is his
quickness to anger (\textsection 37--40). Cicero, who in
later years would write a treatise on the same vice in
the dialogue form (\textit{Tusc.} 4 and the lost
\textit{de Ira}), here lays out the practical advice in
small: when the mind is taken up before reason can
intervene, prepare in advance and meditate that anger
must be resisted; especially restrain the tongue. The
closing paragraphs (\textsection 41--46) raise the
question of fame and posterity (``you do not seek glory
for yourself alone''), liken Asia to a theatre, and end
on the famous closing image: as good poets and
industrious actors are diligent in the last act, so let
Quintus be diligent in this third year of his command,
``that this third year may seem to have been, as the
third act, the most perfect and the most adorned.'' The
final sentence is the older brother's: ``it remains that
I beg you, if you wish me and all yours to be well, that
you most diligently serve your own health.''}
